% "Geometric Dark Energy and the Spin-Torsion Bounce:
%  A Complete Theoretical Assessment"
%
% Houston Golden — 2026
% Preprint: HUBIFY-2026-002
%
% Compile: pdflatex main && pdflatex main (inline bibliography)

\documentclass[aps,prd,twocolumn,superscriptaddress,showpacs,preprintnumbers,
  nofootinbib,longbibliography,floatfix]{revtex4-2}

\usepackage{amsmath,amssymb,amsfonts}
\usepackage{graphicx}
\usepackage{bm}
\usepackage{hyperref}
\usepackage{xcolor}
\usepackage{booktabs}
\usepackage{multirow}
\usepackage{dcolumn}
\usepackage{enumitem}
\usepackage{mathtools}
\usepackage{bbold}
\usepackage[utf8]{inputenc}
\usepackage{float}
\usepackage{slashed}

\hypersetup{
  colorlinks=true,
  linkcolor=blue!60!black,
  citecolor=blue!60!black,
  urlcolor=blue!60!black
}

% Custom commands
\newcommand{\Leff}{\Lambda_{\rm eff}}
\newcommand{\Lobs}{\Lambda_{\rm obs}}
\newcommand{\Dinf}{\mathcal{D}_{\rm inf}}
\newcommand{\rhocrit}{\rho_{\rm crit}}
\newcommand{\rhoPl}{\rho_{\rm Pl}}
\newcommand{\MPl}{M_{\rm Pl}}
\newcommand{\Neff}{N_{\rm eff}}
\newcommand{\BI}{\gamma}
\newcommand{\Geff}{G_{\rm eff}}
\newcommand{\Gsp}{G_{\rm SP}}
\newcommand{\Veff}{V_{\rm eff}}
\newcommand{\OmGW}{\Omega_{\rm GW}}
\newcommand{\mT}{m_{\rm T}}
\newcommand{\geff}{g_{\rm eff}}
\newcommand{\ie}{i.e.}
\newcommand{\eg}{e.g.}
\newcommand{\cf}{cf.}
\newcommand{\etal}{\textit{et~al.}}
\newcommand{\paperVersion}{v2.0.0}

\begin{document}

\preprint{HUBIFY-2026-002}

\title{Geometric Dark Energy and the Spin-Torsion Bounce:\\
A Complete Theoretical Assessment}

\author{Houston Golden}
\email{houston@hubify.com}
\affiliation{Independent Researcher, Los Angeles, California, USA}

\date{\today}

\begin{abstract}
We present a comprehensive theoretical assessment of the Einstein-Cartan-Holst
(ECH) framework as both a dark-energy phenomenology and a bouncing cosmology.
The paper has three parts. First, we describe a phenomenological dark-energy
ansatz rooted in ECH gravity that reduces fine-tuning from $10^{120}$ to
$\sim\!10^5$ and yields MCMC fits partially reducing the $H_0$ and
$\sigma_8/S_8$ tensions ($H_0 = 69.2 \pm 0.8$~km~s$^{-1}$~Mpc$^{-1}$,
$\sigma_8 = 0.785 \pm 0.016$). Second, we report the systematic closure
of four minimal and three extended first-principles routes to deriving
the late-time vacuum term, spanning Foundations A through G (seven
structural barriers). Third, we present a comprehensive characterization
of the spin-torsion bounce's observable prospects across all perturbation
channels: the tensor power spectrum is $P_T \approx 2\times 10^{-64}$
(flat, $n_T \approx 0$), the chirality is $\Delta\chi = 0$ exactly
(the effective four-fermion interaction $(J^5)^2$ is parity-\textit{even}),
the scalar transfer function is $T(k)=1$ for all observable modes
(time-reversal symmetry), and bounce-triggered state selection is
prevented by Liouville's theorem. The Poincar\'e gauge theory (PGT)
extension lowers the bounce scale but the mass-coupling lock
($\geff \sim \mT/\MPl^2$) suppresses torsion-specific features, and the
vacuum amplification ceiling ($\OmGW \propto (H/\MPl)^2$) ensures a
minimum $10^{17}$ gap to any gravitational-wave detector. Two further
routes are tested and closed: baryogenesis and relic production
(Branch~N) fail because at $T\sim\MPl$ the torsion $(J^5)^2$ interaction
is one of $\sim\!100$ gravitational-strength channels (gravitational
democracy), and hidden-sector vacuum selection (Branch~O) fails because
the bounce determines \textit{whether} a transition occurs but not
\textit{what} vacuum energy results (bounce--vacuum energy decoupling).
Finally, the PGT Sector~II pseudoscalar torsion mode---the last candidate
for a torsion-specific relic---is protected by an exact $\mathbb{Z}_2$
parity symmetry that prevents population at the bounce (Branch~P).
Fourteen structural barriers are cataloged across ten research branches
(H through P), spanning five failure modes. The spin-torsion bounce is
theoretically consistent but observationally inert: what is detectable is
generic to any bounce; what is specific to torsion is undetectable.
\end{abstract}

\pacs{98.80.-k, 04.50.Kd, 04.60.Pp, 95.36.+x}
\keywords{dark energy, Einstein-Cartan theory, spin-torsion coupling,
cosmological constant problem, Barbero-Immirzi parameter, bouncing cosmology,
gravitational waves, negative results, Poincar\'e gauge theory}

\maketitle
\tableofcontents


%% ============================================================
%%  1. INTRODUCTION
%% ============================================================
\section{Introduction}\label{sec:intro}

The cosmological constant problem---why the observed vacuum energy density
$\rho_\Lambda \sim (2.3\;\text{meV})^4$ is $\sim 10^{120}$ times smaller
than the naive quantum field theory estimate $\sim \MPl^4$~\cite{Weinberg1989}%
---remains one of the deepest unsolved problems in fundamental physics.
Growing tensions within $\Lambda$CDM, including the ${\sim}4.9\sigma$
Hubble tension~\cite{Riess2022,Planck2018params} and $2$--$3\sigma$
$\sigma_8/S_8$ tension~\cite{KiDS2021,DES2024}, together with hints of
dynamical dark energy from DESI~DR2 BAO
measurements~\cite{DESI2024,DESI2025DR2}, motivate exploration of
geometric approaches to dark energy.

Einstein-Cartan-Holst (ECH) gravity---Einstein-Cartan theory supplemented
by the Holst term with Barbero-Immirzi parameter $\BI$---is a natural
starting point for geometric dark energy. The theory is well-motivated by
loop quantum gravity~\cite{Ashtekar2011}, introduces spin-torsion coupling
with parity-violating structure~\cite{Freidel2005,Mercuri2006}, and
provides a non-singular quantum bounce~\cite{Poplawski2012}. The central
questions are whether this theoretical structure can (i) generate a natural
origin for late-time dark energy, and (ii) produce distinctive observable
signatures from the bounce.

This paper reports three sets of results, organized as a three-part
assessment.

\textit{Part~I: Framework and Phenomenology} (Secs.~\ref{sec:ech}--\ref{sec:ansatz}).
We present the ECH framework and a phenomenological dark-energy ansatz that
is observationally consistent and reduces fine-tuning significantly.
MCMC fits partially reduce the $H_0$ and $\sigma_8$ tensions. The equation
of state $w = -1$ is assumed, not derived.

\textit{Part~II: Dark Energy Derivation Program} (Secs.~\ref{sec:program}--\ref{sec:lessons}).
We report the systematic failure of four minimal first-principles routes to
derive the late-time $w=-1$ behavior, together with the assessment of three
extended foundations (A through G), identifying seven structural barriers
and six structural lessons.

\textit{Part~III: Bounce Observable Signatures} (Secs.~\ref{sec:tensor}--\ref{sec:pgt_gw}).
We present a comprehensive characterization of the spin-torsion bounce
across all observable channels: tensor perturbations (Branch~H), scalar
perturbations (Branch~K), bounce-compatible dark energy (Branch~I), state
selection (Branch~J), UV$\to$IR bridge extensions (Branch~L), and the PGT
gravitational wave spectrum (Branch~M). Five additional barriers are
identified, bringing the total to twelve.

The combination establishes the current scientific status of the program:
\begin{itemize}[leftmargin=*]
  \item The ECH framework is a \textit{motivated phenomenological model},
    not a first-principles derivation of dark energy.
  \item All tested routes to deriving $w=-1$ from geometry have failed,
    for distinct structural reasons spanning seven barriers.
  \item The spin-torsion bounce is observationally inert in the minimal
    model. The PGT extension produces a well-defined but undetectable
    gravitational wave spectrum.
  \item Twelve structural barriers, spanning three failure modes
    (energy-budget, signature, state-selection), prevent any nontrivial
    connection between the bounce and dark energy or between the bounce
    and cosmological observations.
\end{itemize}

\subsection{Relationship to prior work}

This paper supersedes an earlier version~\cite{Golden2026v1} (Paper~1)
that presented the phenomenological framework organized around the minimal
ECH model as if it were a sufficient foundation. The present version
incorporates the results of a structured derivation and bounce
characterization program (Branches H through M) that tested all natural
routes and found all to be negative or generic. The phenomenological
results are unchanged; the framing, scope, and conclusions are
substantially extended to include the bounce sector.


%% ============================================================
%%  2. EINSTEIN-CARTAN-HOLST FRAMEWORK
%% ============================================================
\section{Einstein-Cartan-Holst Framework}\label{sec:ech}

\subsection{Action and field content}\label{sec:action}

The starting point is the Einstein-Cartan-Holst action minimally coupled
to $N_f$ Dirac fermions:
\begin{align}\label{eq:ECH}
  S_{\rm ECH} &= \frac{\MPl^2}{2}\int
    \bigl[\varepsilon_{IJKL} + \tfrac{2}{\BI}\,\eta_{I[K}\eta_{L]J}\bigr]\,
    e^I\wedge e^J\wedge F^{KL}[\omega] \nonumber\\
  &\quad + \int d^4x\,|e|\,\bigl[\bar\psi\,i\gamma^\mu D_\mu\psi
    - m\bar\psi\psi\bigr],
\end{align}
with $D_\mu = \partial_\mu + \tfrac{1}{4}\omega_\mu^{IJ}\gamma_I\gamma_J$
the Lorentz-covariant derivative and $\BI = 0.274 \pm 0.020$ the
Barbero-Immirzi parameter fixed by loop quantum gravity black hole entropy
calculations~\cite{Ashtekar2011}. The first term is the Holst action:
the Palatini formulation of general relativity supplemented by the
parity-odd Holst term $(1/\BI)\eta_{I[K}\eta_{L]J}e^I e^J F^{KL}$.
The second term is the standard minimal Dirac coupling.

\subsection{Torsion equation of motion}\label{sec:torsion}

In ECH gravity, the spin connection $\omega^{IJ}_\mu$ is an independent
dynamical variable, and its equation of motion is algebraic---torsion is
\textit{not propagating}. Solving the torsion equation exactly
decomposes the connection as $\omega = \mathring\omega + K$, where
$\mathring\omega$ is the torsion-free Levi-Civita connection and $K$ is
the contorsion tensor determined algebraically by the fermion spin
density~\cite{Freidel2005,Shapiro2002}.

\subsection{Reduced action after torsion elimination}\label{sec:reduced}

Substituting the algebraic torsion solution back into the action yields
the reduced action~\cite{Freidel2005,ShapiroTeixeira2014}:
\begin{equation}\label{eq:Sreduced}
  S_{\rm red} = S_{\rm EH}[g] + S_{\rm Dirac}[g,\psi]
  + S_{\rm 4f}[\psi;\BI],
\end{equation}
where $S_{\rm EH}$ is the standard Einstein-Hilbert action,
$S_{\rm Dirac}$ is the standard minimally-coupled Dirac action with
the Levi-Civita connection, and $S_{\rm 4f}$ is the torsion-induced
four-fermion contact interaction:
\begin{equation}\label{eq:L4f}
  \mathcal{L}_{\rm 4f}
  = -\Geff\Bigl(\bar\psi\gamma^\mu\gamma^5\psi
    + \tfrac{1}{\BI}\,\bar\psi\gamma^\mu\psi\Bigr)^{\!2},
  \quad
  \Geff = \frac{3\kappa^2}{16}\,\frac{\BI^2}{\BI^2+1}.
\end{equation}
This perfect-square structure constrains the Fierz channel decomposition.
At $\BI = 0.274$, the vector admixture $1/\BI \approx 3.65$ dominates;
the scalar/pseudoscalar channel coupling
$\Gsp \propto (\BI^2 - 1)/(\BI^2 + 1)$ is repulsive (negative).

\paragraph{Critical structural observation.}
After torsion elimination, \textit{all} $\BI$-dependence resides in the
four-fermion term $S_{\rm 4f}$. The Dirac kinetic operator
$i\gamma^\mu\mathring\nabla_\mu$ is the standard Levi-Civita operator,
independent of $\BI$. This structural fact is the root cause of the
failures reported in Sec.~\ref{sec:closures}.

\subsection{Modified Friedmann equation}\label{sec:friedmann}

The spin-torsion four-fermion interaction generates a $\rho^2$ correction
to the Friedmann equation at high densities. When the fermion content is
treated as a perfect fluid with energy density $\rho$, the effective
Friedmann equation reads~\cite{Poplawski2012}:
\begin{equation}\label{eq:modfriedmann}
  H^2 = \frac{8\pi G}{3}\,\rho\Bigl(1 - \frac{\rho}{\rhocrit}\Bigr),
\end{equation}
with critical density
\begin{equation}\label{eq:rhocrit}
  \rhocrit \simeq 0.21\,\MPl^4,
\end{equation}
where the numerical prefactor depends on the number of fermion species and
$\BI$. Equation~\eqref{eq:modfriedmann} implies a non-singular bounce:
$H=0$ when $\rho = \rhocrit$, and $\dot{a} > 0$ for $\rho < \rhocrit$.
The bounce replaces the classical Big Bang singularity with a smooth
transition from contraction to expansion.


%% ============================================================
%%  3. PHENOMENOLOGICAL DARK-ENERGY ANSATZ
%% ============================================================
\section{Phenomenological Dark-Energy Ansatz}\label{sec:ansatz}

\subsection{The scaling argument}\label{sec:scaling}

The ECH framework contains a parity-odd vacuum operator
$\varepsilon^{abcd}K_{ab}R_{cd}$ sourced by the fermion spin density via
torsion. At early times (Planck-density bounce), this operator has a
natural magnitude set by $\alpha/M \sim \mathcal{O}(\kappa^2)$, where
$\alpha/M$ is the parity-odd coefficient in the effective Lagrangian.

Inflationary expansion after the bounce dilutes this operator by a factor
$\Dinf \sim e^{-3N_{\rm tot}}$, where $N_{\rm tot} \sim 92$ $e$-folds.
The resulting late-time dark-energy density is parametrically
\begin{equation}\label{eq:rhoL}
  \rho_\Lambda \approx \Xi\,\MPl^4, \qquad
  \Xi \equiv \bigl[(\alpha/M)\,\MPl\bigr]\,\Dinf,
\end{equation}
giving $\Xi \sim 10^{-122}$ for appropriate values of $\alpha/M$ and
$N_{\rm tot}$.

\paragraph{What this is.}
Equation~\eqref{eq:rhoL} is a \textit{phenomenological scaling ansatz}.
It traces a dimensional chain from the Planck scale through inflationary
dilution to the observed vacuum energy. It reduces fine-tuning from
$10^{120}$ (the cosmological constant problem in $\Lambda$CDM) to
$\sim\!10^5$ (sensitivity to $N_{\rm tot}$ within $\Delta N \approx 4$
$e$-folds).

\paragraph{What this is not.}
Equation~\eqref{eq:rhoL} is \textit{not} a derivation of $w = -1$.
The parity-odd operator is sourced by the spin density, which vanishes
at late times ($K_{ab} \to 0$). For the residual to persist as a true
vacuum energy, one must show that integrating out the spin degrees of
freedom generates an IR-constant term in the effective action. Four
attempts to establish this have failed (Sec.~\ref{sec:closures}). We
therefore treat $w = -1$ as a phenomenological assumption, consistent
with observations ($|w+1| \lesssim 10^{-2}$ from DESI+Planck).

\subsection{Cosmological parameters}\label{sec:cosmo}

The phenomenological framework motivates a $\Lambda$CDM$+\Delta\Neff$
extension: the bounce scenario qualitatively predicts additional
relativistic species from particle production near the Planck density,
but quantitative estimates are not yet available.

\subsubsection{MCMC results}

MCMC fits to the full tension dataset (Planck~2018 + DESI~DR1 BAO +
Pantheon+ + SH0ES $H_0$ prior + DES~Y3 $S_8$) yield:
\begin{align}
  H_0 &= 69.2 \pm 0.8 \;\text{km}\;\text{s}^{-1}\;\text{Mpc}^{-1},
    \label{eq:H0}\\
  \sigma_8 &= 0.785 \pm 0.016, \label{eq:s8}\\
  S_8 &= 0.80 \pm 0.02. \label{eq:S8}
\end{align}
These partially reduce the $H_0$ tension (from $4.9\sigma$ to
$2.9\sigma$ vs.\ SH0ES) and the $S_8$ tension (from $\sim\!3\sigma$ to
$\sim\!1\sigma$ vs.\ weak lensing).

\begin{table}[t]
\caption{\label{tab:mcmc} Cosmological parameters from MCMC fits.
Posterior means $\pm\,1\sigma$. Original analysis uses the full tension
dataset; verification uses Cobaya~v3.6.1 with CamSpec.}
\begin{ruledtabular}
\begin{tabular}{lcc}
\textbf{Parameter} & \textbf{Original} & \textbf{Verification} \\
\hline
$H_0$ [km/s/Mpc] & $69.2 \pm 0.8$ & $67.68 \pm 1.06$ \\
$\sigma_8$ & $0.785 \pm 0.016$ & $0.803 \pm 0.008$ \\
$S_8$ & $0.80 \pm 0.02$ & $0.814 \pm 0.008$ \\
$\Delta\Neff$ & --- & $-0.020 \pm 0.169$ \\
$\Omega_m$ & $0.310 \pm 0.008$ & $0.308 \pm 0.005$ \\
$n_s$ & $0.9649 \pm 0.0042$ & --- \\
$\tau$ & $0.054 \pm 0.007$ & --- \\
\hline
Total samples & --- & 176{,}840 \\
$\hat{R}-1$ & $< 0.01$ & $0.001$ \\
\end{tabular}
\end{ruledtabular}
\end{table}

\subsubsection{Independent verification}\label{sec:verification}

Independent verification using Cobaya~v3.6.1 with Planck NPIPE CamSpec
TTTEEE + lowl TT/EE + lensing finds $\Delta\Neff$ consistent with zero
($\Delta\Neff = -0.020 \pm 0.169$ for the full-tension dataset,
$\Delta\Neff = +0.065 \pm 0.17$ for Planck+BAO+SN without $H_0$/$S_8$
priors; see Table~\ref{tab:mcmc}). The tension reduction in the
original analysis is driven by the SH0ES prior, not by the $\Delta\Neff$
extension itself.

\subsection{Fine-tuning reduction}

A Monte Carlo sensitivity scan ($10^5$ samples) confirms that the
scaling ansatz Eq.~\eqref{eq:rhoL} reduces fine-tuning from $10^{120}$
to $\sim\!10^5$: the viable range $N_{\rm tot} \in [79,\,95]$ produces
$\rho_\Lambda$ within two orders of magnitude of the observed value,
with $N_{\rm tot}$ as the single controlling parameter ($|\rho_s| =
0.996$). The residual tuning is concentrated in a single initial
condition with potential dynamical answers through the
bounce-to-inflation transition.


%% ============================================================
%%  PART II: SYSTEMATIC MODEL TESTING
%% ============================================================

%% ============================================================
%%  4. THE DARK ENERGY DERIVATION PROGRAM
%% ============================================================
\section{Dark Energy Derivation Program}\label{sec:program}

The central question is whether late-time $w = -1$ can be derived from
first principles within the minimal ECH$+$Dirac model or its natural
extensions. We tested four minimal routes and assessed three extended
foundations (A through G), each with predefined gates and kill criteria
frozen before computation.

\subsection{Four minimal routes tested}\label{sec:routes}

\begin{description}[style=nextline,leftmargin=1em]
  \item[Track B (NJL condensate)]
    Test whether a pseudoscalar fermion condensate
    $\langle\bar\psi i\gamma^5\psi\rangle$ forms via the
    Nambu-Jona-Lasinio mechanism from the torsion-induced four-fermion
    interaction and persists as late-time vacuum energy.

  \item[Branch G v1 (strict one-loop effective action)]
    Compute the one-loop gravitational effective action $\Gamma[e]$ after
    torsion elimination and fermion integration, and test for
    $\BI$-dependent vacuum content.

  \item[Route T1 (dynamical Immirzi field)]
    Promote $\BI \to \theta(x)$ to a dynamical pseudoscalar and assess
    whether the reduced action contains novel content not reproducible
    by adding a standard axion-like particle (ALP) to GR.

  \item[Route S1 (parity-sensitive CMB phenomenology)]
    Assess whether the framework's parity-odd operator maps to distinctive
    birefringence observables distinguishable from generic ALP predictions.
\end{description}

All four routes were closed with clean negative results. Additionally,
Foundations~A through~G (an extended program testing propagating torsion,
Nieh-Yan non-topological coupling, environmental mass, disformal coupling,
scale separation, attractor sensitivity, and vacuum selection) were all
closed, yielding seven barriers in total.


%% ============================================================
%%  5. CLOSURE RESULTS
%% ============================================================
\section{Closure Results}\label{sec:closures}

\subsection{Track B: condensate route}\label{sec:trackB}

Fierz rearrangement of the four-fermion
interaction~\eqref{eq:L4f} decomposes it into standard NJL channels.
The scalar/pseudoscalar coupling is
\begin{equation}\label{eq:GSP}
  \Gsp = \frac{3\kappa^2}{16}\,\frac{\BI^2 - 1}{\BI^2 + 1},
\end{equation}
which is \textit{repulsive} (negative) for $\BI < 1$. At $\BI = 0.274$,
the NJL gap equation has no nontrivial solution. Even for $\BI > 1$
(attractive sign), the coupling is $\sim\!175\times$ below the NJL
critical coupling. Track~B closes at Gate~1: no condensate forms.

\subsection{Branch G v1: strict one-loop}\label{sec:branchG}

After torsion elimination, the Dirac operator is
$\slashed{D} = i\gamma^\mu\mathring\nabla_\mu$---the standard
$\BI$-independent Levi-Civita operator. The four-fermion term is quartic
in $\psi$ and does not contribute at one loop. The Seeley-DeWitt
coefficients $a_0, a_1, a_2$ are the standard textbook values with no
$\BI$-dependence~\cite{Vassilevich2003}. Branch~G v1 closes: no novel
Holst content at strict one-loop.

\subsection{Route T1: dynamical Immirzi field}\label{sec:T1}

Promoting $\BI \to \theta(x)$, torsion remains algebraic. After
elimination, the reduced action is~\cite{CalcagniMercuri2009,Mercuri2009}
\begin{equation}\label{eq:T1red}
  S_{\rm red} = S_{\rm EH} + S_{\rm Dirac}
  + \int\!\bigl[-\tfrac{1}{2}f(\theta)(\partial\theta)^2
    + \mathcal{L}_{\rm 4f}\bigr],
\end{equation}
where $\mathcal{L}_{\rm 4f}$ is $\theta$-independent. Every term is
reproducible by adding a standard ALP to GR. The Nieh-Yan 4-form is
exact (total derivative) in 4D; its topological nature ensures the
gravitational origin washes out~\cite{CalcagniMercuri2009}. Cosmological
dynamics give $w = +1$ (stiff matter), the opposite of dark
energy~\cite{TaverasYunes2008}. Route~T1 closes.

\subsection{Route S1: parity CMB phenomenology}\label{sec:S1}

The minimal ECH action has no photon-polarization coupling. The
parity-odd operator is purely gravitational; the electromagnetic field
strength does not appear in the Holst term or the reduced action. Any
assumed coupling produces constant-$\beta$ birefringence with two free
parameters for one observable---identical to generic ALP phenomenology.
Route~S1 closes.

\subsection{Foundations A--G: extended program}\label{sec:foundations}

Beyond the four minimal routes, three extended foundations were assessed:
\begin{itemize}[leftmargin=*]
  \item \textit{Foundation~A} (PGT propagating torsion): The ghost-free
    $0^-$ axial torsion mode has mass $m_B = \MPl/(4\sqrt{\pi|t_3|})$
    and matter coupling $\geff \sim 1/(\MPl\sqrt{|t_3|})$. The
    \textit{mass-coupling lock} (Barrier~1): $m_B$ and $\geff$ decrease
    together, with $m_B/\geff \sim \MPl^2$ fixed. At meV-scale mass,
    the coupling is $10^{29}\times$ below gravitational strength. No
    symmetry within PGT protects $m_B$.
  \item \textit{Foundations~B--D}: The Nieh-Yan non-topological route
    (B), environmental mass (C), and disformal coupling (D) were all
    closed by distinct mechanisms: topological-shift duality, scalar-tensor
    universality, and Planck suppression, respectively.
  \item \textit{Foundations~E--G}: Scale separation, attractor-sensitivity
    dilemma, and parameter immunity closed the remaining directions.
\end{itemize}

\subsection{Closure summary}\label{sec:closure_summary}

\begin{table}[t]
\centering
\caption{\label{tab:closures} Tested routes and outcomes. Four minimal
routes plus Foundation~A (the most natural PGT extension).}
\begin{ruledtabular}
\begin{tabular}{llll}
\textbf{Route} & \textbf{Target} & \textbf{Method} & \textbf{Failure} \\
\hline
Track B & Condensate & Fierz + NJL & $\Gsp < 0$ at $\BI < 1$ \\
G v1 & $\Gamma[e]$ & 1-loop & $\slashed{D}$ is $\BI$-free \\
T1 & $\theta(x)$ & Lit.\ review & $=$ standard ALP \\
S1 & Birefr. & Assessment & No photon coupling \\
Found.\ A & PGT $0^-$ & Phase 1+2 & Mass-coupling lock \\
Found.\ B & Nieh-Yan & MAG analysis & Top.-shift duality \\
Found.\ C & Environ. & Scalar-tensor & FRW universality \\
Found.\ D & Disformal & Connection & Planck suppression \\
Found.\ E & Integrals & Scale sep. & $k_{\rm CMB}/k_b$ \\
Found.\ F & Attractors & IC analysis & Sensitivity dilemma \\
Found.\ G & Vacuum sel. & Cyclic & Parameter immunity \\
\end{tabular}
\end{ruledtabular}
\end{table}


%% ============================================================
%%  6. STRUCTURAL LESSONS FROM THE DE PROGRAM
%% ============================================================
\section{Structural Lessons}\label{sec:lessons}

The closures reveal a hierarchy of structural obstructions. Lessons~1--5
emerge from the minimal ECH model; Lesson~6 from the PGT extension.

\subsection{Lesson 1: Algebraic torsion washes out}

In the minimal ECH model, torsion is non-propagating. Its equation of
motion is algebraic: contorsion $K$ is determined pointwise by the fermion
spin density. Torsion elimination is exact and complete---no torsion-sector
information survives into the reduced action except through the
four-fermion coupling constants.

\subsection{Lesson 2: Fixed $\BI$ leaves no IR fingerprint}

$\BI$ enters the reduced action only through $\Geff$ and $\Gsp$ in the
four-fermion interaction. These are UV-scale contact couplings. At one
loop, the fermion determinant is $\BI$-independent. At the condensate
level, the coupling is subcritical. $\BI$ sets the UV-scale spin-torsion
structure but does not generate low-energy phenomena.

\subsection{Lesson 3: Dynamical $\BI$ reduces to generic ALP}

The Nieh-Yan 4-form $N_4 = d(e^I \wedge T_I)$ is exact in 4D. Any
coupling of a pseudoscalar $\theta$ to $N_4$ reduces to boundary terms
plus standard kinetic/coupling structures. The gravitational origin is
topologically protected from leaving a local dynamical imprint.

\subsection{Lesson 4: Parity alone is not a mechanism}

Without a derived coupling to observable sectors, the ECH action's
parity-odd sector produces no testable prediction distinguishing the
framework from generic alternatives. Parity is a structural property,
not by itself a mechanism or signal source.

\subsection{Lesson 5: No scale protection in the minimal model}

None of the tested routes addresses naturalness. The condensate route
produces the wrong sign. The one-loop route has standard renormalization.
The dynamical Immirzi route requires $V(\theta_0) = \Lobs$ by hand.

\subsection{Lesson 6: The mass-coupling lock}\label{sec:lesson6}

Even when torsion propagates (PGT), the parameter that controls the
torsion mass also controls the matter coupling:
\begin{equation}\label{eq:lock}
  m_B \propto \frac{1}{\sqrt{|t_3|}}, \qquad
  \geff \propto \frac{1}{\sqrt{|t_3|}}, \qquad
  \frac{m_B}{\geff} \sim \MPl^2.
\end{equation}
Making the mode cosmologically light automatically makes it
cosmologically inert. The lock occurs generically when a single parameter
controls both the kinetic normalization and the mass of a geometric mode.
No symmetry within PGT protects the mass, and graviton loops generate
$\delta m_B^2 \sim \MPl^2/(16\pi^2)$, requiring fine-tuning of $1$ in
$10^{57}$ for meV-scale mass.


%% ============================================================
%%  PART III: BOUNCE OBSERVABLE SIGNATURES
%% ============================================================

%% ============================================================
%%  7. BOUNCE TENSOR SPECTRUM (Branch H)
%% ============================================================
\section{Bounce Tensor Spectrum}\label{sec:tensor}

\subsection{Background solution}\label{sec:background}

For a radiation-dominated fluid ($p = \rho/3$) on the modified Friedmann
equation~\eqref{eq:modfriedmann}, the exact scale factor is
\begin{equation}\label{eq:scale_factor}
  a(t) = a_b\bigl(1 + 4\alpha^2 t^2\bigr)^{1/4},
\end{equation}
where $a_b$ is the minimum scale factor at the bounce ($t=0$) and
\begin{equation}\label{eq:alpha}
  \alpha = \sqrt{\frac{8\pi}{3}}\,\MPl
\end{equation}
is set by the Planck scale. At late times ($\alpha t \gg 1$),
Eq.~\eqref{eq:scale_factor} reduces to $a \propto t^{1/2}$, the standard
radiation-dominated expansion.

\subsection{Tensor perturbation equation}

Tensor perturbations $h_{ij}$ on the bounce background satisfy the
standard wave equation in conformal time $\eta$:
\begin{equation}\label{eq:tensor_mode}
  \mu_k'' + \Bigl(k^2 - \frac{a''}{a}\Bigr)\mu_k = 0,
\end{equation}
where $\mu_k = a\,h_k$ is the Mukhanov variable, primes denote
derivatives with respect to conformal time, and the effective potential is
\begin{equation}\label{eq:tensor_potential}
  \frac{a''}{a} = \frac{2\alpha^2 a_b^2}{\bigl(1 + 4\alpha^2 t^2\bigr)^{3/2}}.
\end{equation}
The potential is \textit{always positive}, localized within a time
$\Delta t \sim 1/\alpha \sim t_{\rm Pl}$ around the bounce, and defines a
characteristic wavenumber $k_b = a_b\sqrt{2}\,\alpha \approx
1.88\,a_b\,\MPl$.

\subsection{Bogoliubov coefficients}

Numerical solution of Eq.~\eqref{eq:tensor_mode} with vacuum initial
conditions in the contracting phase yields the Bogoliubov coefficient:
\begin{equation}\label{eq:bogoliubov}
  |\beta_k|^2 \approx 0.855\,\frac{k_b^2}{k^2}
  \quad \text{for } k \ll k_b,
\end{equation}
with exponential suppression for $k \gg k_b$. The $1/k^2$ scaling is
confirmed analytically from the Born approximation for scattering off
the localized potential.

\subsection{Tensor power spectrum}

The tensor power spectrum is
\begin{equation}\label{eq:PT}
  P_T(k) \propto k^2\bigl(1 + 2|\beta_k|^2\bigr) \approx \text{const}
  \quad \text{for } k \ll k_b.
\end{equation}
The $k^2$ phase-space factor exactly cancels the $1/k^2$ from the
Bogoliubov coefficient, producing a \textit{flat} (scale-invariant)
spectrum with
\begin{equation}\label{eq:nT}
  n_T \approx 0 \quad (\text{numerical: } n_T = -0.002 \pm 0.01).
\end{equation}
The absolute amplitude, however, is suppressed by the enormous redshift
from the bounce epoch to the present:
\begin{equation}\label{eq:PT_amplitude}
  \boxed{P_T \approx 2 \times 10^{-64},}
\end{equation}
reflecting a dilution factor $(a_b/a_0)^2 \sim 10^{-65}$. The
tensor-to-scalar ratio is correspondingly
\begin{equation}\label{eq:r}
  r = \frac{P_T}{P_S} \sim \frac{10^{-64}}{10^{-9}} \sim 10^{-55},
\end{equation}
which is 51 orders of magnitude below any conceivable detection
threshold.

\subsection{Chirality analysis}

A key question is whether the spin-torsion bounce produces chiral
gravitational waves (different amplitudes for left- and right-circular
polarizations). Six candidate parity-odd mechanisms were analyzed:

\begin{table}[t]
\centering
\caption{\label{tab:parity} Parity-odd candidate mechanisms. None
provides a tensor chirality source within the minimal model.}
\begin{ruledtabular}
\begin{tabular}{lcc}
\textbf{Candidate} & \textbf{Parity-odd?} & \textbf{Tensor coupling?} \\
\hline
$(J^5)^2$ interaction & No (even) & --- \\
$n_5$ background & Conditional & No \\
Holst/Immirzi & Rescales only & No \\
Dynamical Nieh-Yan & Yes (not EC) & Generic ALP \\
Chern-Simons $\tilde{R}R$ & Yes (not EC) & Yes (not EC) \\
Chiral anomaly & Yes (quantum) & No ($\tilde{R}R=0$ on FRW) \\
\end{tabular}
\end{ruledtabular}
\end{table}

The result is unambiguous:
\begin{equation}\label{eq:chirality}
  \boxed{\Delta\chi(k) \equiv \frac{P_L - P_R}{P_L + P_R} = 0
  \quad \text{(exact)}.}
\end{equation}
The four-fermion interaction $(J^5_\mu)^2$ is the product of two
pseudovectors, giving a \textit{scalar}---it is parity-\textit{even}
despite arising from parity-odd torsion. FRW isotropy eliminates all
spatial parity-odd backgrounds ($\langle J^5_i\rangle = 0$), and the
Pontryagin density vanishes identically on FRW
($\tilde{R}^{\mu\nu}_{\;\;\;\rho\sigma}R^{\rho\sigma}_{\;\;\;\mu\nu}=0$).
The bounce breaks time-reversal symmetry, not spatial parity.

This constitutes \textit{Barrier~8: parity-even effective interaction}.
The spin-torsion four-fermion interaction is parity-even; FRW isotropy
eliminates all remaining parity-odd backgrounds; no chirality in tensor
observables.

\subsection{Characteristic frequency}

The bounce tensor spectrum extends from $f=0$ to a cutoff at
\begin{equation}\label{eq:fb_minimal}
  f_b \approx 8\;\text{GHz}
\end{equation}
in the minimal model. This is far above any gravitational wave detector
band.


%% ============================================================
%%  8. SCALAR PERTURBATIONS (Branch K)
%% ============================================================
\section{Scalar Perturbations}\label{sec:scalar}

\subsection{Bardeen equation through the bounce}

Scalar perturbations of a radiation fluid on the bounce background are
described by the Bardeen potential $\Phi$, which satisfies
\begin{equation}\label{eq:bardeen}
  \ddot\Phi + 5H\dot\Phi +
  \Bigl[\frac{k^2}{3a^2} + 2\dot H + 4H^2\Bigr]\Phi = 0.
\end{equation}
Crucially, this equation is \textit{regular} at $H=0$ (the bounce
instant), unlike the Mukhanov-Sasaki variable $v = z\mathcal{R}$ where
$z \propto a\dot\phi/H$ diverges at the bounce. The Bardeen formulation
avoids this singularity entirely, as noted by Peter and
Pinto-Neto~\cite{PeterPintoNeto2008}.

\subsection{Transfer function}

\paragraph{Analytic proof.}
The Bardeen equation possesses exact time-reversal symmetry on the
spin-torsion bounce background. Under $\eta \to -\eta$, the growing mode
of contraction maps to the decaying mode of expansion, and the constant
mode maps to itself with unit amplitude. Therefore
\begin{equation}\label{eq:transfer}
  \boxed{T(k) = 1 \quad \text{for all } k \ll k_b.}
\end{equation}
The bounce does not modify the scalar power spectrum at any observable
scale. Features appear only at $k \sim k_b$ (GHz frequencies).

\paragraph{Numerical confirmation.}
$T(k) = 1.000$ to machine precision for $k/k_b < 0.1$. The growing mode
amplifies by $(k_b/k)^3 \sim 10^{84}$ for CMB modes during contraction,
but the time-reversal symmetry maps this entirely to the decaying mode
of the expanding phase. No matching ambiguity, no spectral distortion.

\subsection{Non-Gaussianity}

The torsion-specific $(J^5)^2$ interaction generates
\begin{equation}\label{eq:fNL}
  f_{\rm NL}^{\rm torsion} \sim 10^{-56}
\end{equation}
at CMB scales, suppressed by the $(k/k_b)^2$ localization of the bounce
interaction. Generic mode coupling through gravity gives $f_{\rm NL}
\sim \mathcal{O}(1)$ at most---the same as any bounce. No distinctive
non-Gaussian signal.

\subsection{Pre-bounce dependence}

The observed scalar spectrum ($A_s$, $n_s$, running) is entirely
determined by the pre-bounce contraction mechanism, which the minimal
model does not specify. A pure radiation contraction gives
$P_S \propto k^4$ (excluded by CMB data). A viable spectrum requires
a matter-dominated or ekpyrotic pre-bounce phase---physics beyond the
minimal spin-torsion bounce.

\subsection{Summary: scalar sector}

The scalar sector provides \textit{no} distinctive bounce signature at
observable scales. The scale separation $k_{\rm CMB}/k_b \sim 10^{-28}$
ensures that the bounce is a sub-resolution event for all cosmological
probes. Everything observable (the spectral shape, the transfer function)
is generic to any time-symmetric radiation bounce; everything specific to
spin-torsion (the bounce profile, the $(J^5)^2$ correction) is confined
to GHz scales.


%% ============================================================
%%  9. BOUNCE-COMPATIBLE DARK ENERGY AND STATE SELECTION
%% ============================================================
\section{Bounce-Compatible Dark Energy and State Selection}\label{sec:compat}

\subsection{Horndeski stability at the bounce (Branch I)}\label{sec:horndeski}

We tested whether the spin-torsion bounce ($H=0$, $\dot{H} \approx
3.52\,\MPl^2$) constrains which Horndeski/scalar-tensor dark energy
models are viable. Six model classes were screened for ghost freedom,
gradient stability, and EFT validity at the bounce:

\begin{table}[t]
\centering
\caption{\label{tab:horndeski} Horndeski model classes at the bounce.}
\begin{ruledtabular}
\begin{tabular}{lll}
\textbf{Class} & \textbf{Verdict} & \textbf{Key finding} \\
\hline
Quintessence & Trivial compat. & $(Q_S,c_S^2) = (1,1)$ \\
K-essence & Trivial compat. & $X \to 0$ at bounce \\
Braiding & EFT breaks & $Q_S \to 0$; $M \sim H_0 \ll \MPl$ \\
Non-minimal & Compatible & Tachyonic kick, bounded \\
Quartic/Quintic & Trivial compat. & Pre-empted by GW170817 \\
DHOST & EFT breaks & Degeneracy algebraic \\
\end{tabular}
\end{ruledtabular}
\end{table}

The root cause is scale separation: $\rho_{\rm DE}/\rhocrit \sim
10^{-122}$. DE fields are frozen ($X \to 0$, $\dot\phi \sim 10^{-60}\,
\MPl^2$) during the Planck-scale bounce. All stability conditions reduce
to their trivial limits.

A methodological finding: the EFT-of-DE $\alpha$-parameterization
(Bellini-Sawicki) is singular at $H=0$ ($\alpha_K \propto 1/H^2$,
$\alpha_B \propto 1/H$ diverge), while physical stability quantities
remain finite. Standard Boltzmann codes cannot be applied to bouncing
cosmologies without reformulation.

\subsection{State selection (Branch J)}\label{sec:state}

Five candidate mechanisms for bounce-triggered dark-energy state selection
were tested: pNGB misalignment, multi-vacuum selection, symmetry
re-breaking, metastable trapping, and nonadiabatic excitation. All five
fail for three converging reasons:

\paragraph{Barrier 9: Hamiltonian phase-space conservation.}
The bounce is a Hamiltonian scattering event occurring near $H=0$
(negligible friction). Liouville's theorem guarantees that the mapping
from pre-bounce to post-bounce states preserves phase-space volume. The
bounce \textit{rotates} dark-sector initial conditions but cannot
\textit{contract} them to a predictive outcome. Phase-space contraction
requires dissipation or irreversibility, both absent at the bounce
instant.

\paragraph{Naturalness dilemma.}
Any curvature coupling strong enough to affect the dark sector at the
bounce ($\xi \sim \mathcal{O}(1)$) generates Planck-scale radiative
corrections to the DE mass ($\delta m^2 \sim \xi\,\MPl^2$), destroying
the DE potential structure. A coupling weak enough to preserve the DE
sector ($\xi \sim 10^{-124}$) is far too weak for state selection.

\paragraph{Scale separation.}
The $10^{61}$ hierarchy between $\MPl$ and $m_{\rm DE}$ prevents
parametric resonance (frequency mismatch), meaningful field evolution
during the bounce ($t_{\rm Pl} \ll 1/m_{\rm DE}$), and nonadiabatic
effects at DE frequencies.


%% ============================================================
%%  10. UV->IR BRIDGE EXTENSIONS (Branch L)
%% ============================================================
\section{UV$\to$IR Bridge Extensions}\label{sec:bridge}

The scale separation $k_{\rm CMB}/k_b \sim 10^{-28}$ is the meta-barrier
underlying many of the null results above. Branch~L screened seven
candidate mechanisms for transferring bounce-scale physics to observable
scales.

\subsection{Screening results}

\begin{itemize}[leftmargin=*]
  \item \textit{PGT lower-scale bounce}: Survives. In Poincar\'e gauge
    theory, torsion propagates with mass $\mT$, and the effective bounce
    critical density becomes
    \begin{equation}\label{eq:rhocrit_pgt}
      \rhocrit^{\rm eff} \sim \mT^2\,\MPl^2.
    \end{equation}
    For $\mT \sim 10^3$--$10^9$~GeV, bounce features move to the
    LISA--ET frequency bands.
  \item \textit{Bounce + brief inflation}: Killed---reduces to standard
    inflation with the bounce irrelevant.
  \item \textit{Time-asymmetric matter bounce}: Killed---not
    bounce-specific (any bounce works).
  \item \textit{Generic curvaton}: Killed---not bounce-specific.
  \item \textit{Torsion-curvaton}: Conditional ($\sim\!3\%$)---requires
    PGT with specific parameter tuning.
  \item \textit{Bounce-triggered relics}: Killed---not bounce-specific.
  \item \textit{Cyclic resonance}: Killed---pathological (Foundation~G
    incompatibility).
\end{itemize}

The screening reveals a sharp dichotomy (\textit{Barrier~10: UV$\to$IR
specificity dilemma}): mechanisms that beat scale separation generically
(contraction-era spectrum generation, inflation) do not need the bounce,
while the only mechanism that is bounce-specific (lowering $\rhocrit$ via
PGT) requires changing the gravitational theory itself.

\subsection{PGT bounce parameters}

The PGT bounce frequency scales as
\begin{equation}\label{eq:fb_pgt}
  f_b \approx 2.6 \times 10^{10}\,
  \Bigl(\frac{\mT}{\MPl}\Bigr)^{\!1/2}\;\text{Hz},
\end{equation}
and the torsion mass determines the bounce scale:
\begin{equation}\label{eq:mT}
  \mT = \frac{\MPl}{2\sqrt{|t_3|}}.
\end{equation}

\begin{table}[t]
\centering
\caption{\label{tab:pgt_scan} PGT bounce parameter scan: torsion mass
$\mT$ vs.\ bounce frequency $f_b$ and peak GW amplitude
$\OmGW^{\rm peak}h^2$. The detector gap is the ratio of detector
sensitivity to signal amplitude.}
\begin{ruledtabular}
\begin{tabular}{lccc}
$\mT$ [GeV] & $f_b$ [Hz] & $\OmGW^{\rm peak}h^2$ & Gap to best detector \\
\hline
$10^{18}$ ($=\MPl$) & $2.6\times 10^{10}$ & $5\times 10^{-6}$ & $10^4$ (HF) \\
$10^{14}$ & $2.6\times 10^{8}$ & $5\times 10^{-14}$ & $10^9$ (ET) \\
$10^{10}$ & $2.6\times 10^{6}$ & $5\times 10^{-22}$ & $10^{12}$ (ET) \\
$10^{7}$ & $2.4\times 10^{4}$ & $3\times 10^{-30}$ & $10^{17}$ (ET) \\
$10^{4}$ & $2.6\times 10^{2}$ & $5\times 10^{-34}$ & $10^{21}$ (LIGO) \\
$10^{0}$ & $2.6$ & $5\times 10^{-42}$ & $10^{32}$ (DECIGO) \\
$10^{-3}$ & $8.2\times 10^{-2}$ & $5\times 10^{-48}$ & $10^{35}$ (LISA) \\
\end{tabular}
\end{ruledtabular}
\end{table}


%% ============================================================
%%  11. PGT GHOST-FREE PARAMETER SPACE (Branch L Phase 2)
%% ============================================================
\section{PGT Ghost-Free Parameter Space}\label{sec:ghost}

\subsection{Sector II: spin-$0^-$ axial torsion}

In the quadratic PGT action, the ghost-free parameter subspace has been
mapped by Sezgin and van
Nieuwenhuizen~\cite{SezginNieuwenhuizen1980} and
Karananas~\cite{Karananas2015}. The cleanest single-mode model is the
spin-$0^-$ (pseudoscalar axial torsion), which is ghost-free for all
$|t_3| > 0$ and has one light mode.

\subsection{Sector I: rejected}

The spin-$0^+$ scalar torsion mode develops a ghost at the target mass
scale required for observable bounce features. This sector is rejected.

\subsection{Mass-coupling lock reappears}

In the ghost-free $0^-$ sector, canonical normalization gives
\begin{equation}\label{eq:geff_pgt}
  \geff \sim \frac{\mT}{\MPl^2},
\end{equation}
confirming the mass-coupling lock (Lesson~6): the effective matter
coupling is $\mT/\MPl^2$, which vanishes in the light-mass limit.

This constitutes \textit{Barrier~11: decoupling universality}. The
pattern is universal: any geometric mode whose kinetic term is inherited
from the gravitational action without an independent mass-generating
mechanism will exhibit this decoupling in the canonically normalized
theory.


%% ============================================================
%%  12. PGT BOUNCE GRAVITATIONAL WAVE SPECTRUM (Branch M)
%% ============================================================
\section{PGT Bounce Gravitational Wave Spectrum}\label{sec:pgt_gw}

\subsection{Background evolution}

The PGT bounce background has the same functional form as the minimal
model, Eq.~\eqref{eq:scale_factor}, but with $\alpha$ replaced by
\begin{equation}\label{eq:alpha_pgt}
  \alpha_{\rm PGT} = \mT\sqrt{\frac{8\pi}{3}},
\end{equation}
reflecting the lowered bounce energy scale $\rhocrit^{\rm eff} \sim
\mT^2\,\MPl^2$.

\subsection{Effective potential and mode equation}

The tensor mode equation is identical to Eq.~\eqref{eq:tensor_mode},
with the effective potential
\begin{equation}\label{eq:Veff_pgt}
  \frac{a''}{a} = \frac{2\alpha_{\rm PGT}^2\,a_b^2}
  {\bigl(1 + 4\alpha_{\rm PGT}^2\,t^2\bigr)^{3/2}}.
\end{equation}

\subsection{Bogoliubov spectrum}

The Bogoliubov coefficient follows the same functional form as the
minimal model:
\begin{equation}\label{eq:beta_pgt}
  |\beta_k|^2 \approx
  \begin{cases}
    \text{const} & k \ll k_b, \\
    \exp(-\pi k/k_b) & k \gg k_b.
  \end{cases}
\end{equation}

\subsection{GW energy density: the critical result}

The GW energy density spectrum is
\begin{equation}\label{eq:OmGW}
  \boxed{\OmGW(f)\,h^2 \approx 5\times 10^{-6}\,
  \Bigl(\frac{\mT}{\MPl}\Bigr)^{\!2}\,S(f/f_b),}
\end{equation}
where $S(x)$ is a universal shape function with $S \approx 1$ for $x \ll
1$, oscillatory with $1/x^2$ envelope near $x \sim 1$, and exponentially
suppressed for $x \gg 1$.

The critical feature is the factor $(\mT/\MPl)^2$---this is \textit{not}
$\mathcal{O}(1)$ as might be naively expected. The efficiency of vacuum
amplification is
\begin{equation}\label{eq:efficiency}
  \varepsilon = \frac{\rho_{\rm GW}}{\rhocrit} \sim
  \frac{\mT^2}{\MPl^2} = \Bigl(\frac{\mT}{\MPl}\Bigr)^{\!2},
\end{equation}
which is the standard graviton production efficiency: $\mathcal{O}(1)$
gravitons per mode, but each graviton carries energy $\sim \mT$ while
the total energy budget is $\sim \mT\,\MPl$. The ratio is $(\mT/\MPl)^2$.

\subsection{Amplitude-frequency tradeoff}

The peak amplitude satisfies
\begin{equation}\label{eq:tradeoff}
  \OmGW^{\rm peak} \propto f_b^4,
\end{equation}
where $f_b \propto \mT^{1/2}$. Lowering $f_b$ by a factor of $10$
reduces the amplitude by $10^4$. Moving from $f_{\rm Pl} \sim 10^{10}$~Hz
to LISA frequencies ($10^{-3}$~Hz) requires $10^{13}$ in frequency,
giving $10^{52}$ in amplitude suppression. This tradeoff is universal to
\textit{all} vacuum-amplification bounces.

\subsection{Minimum detector gap}

\begin{table}[t]
\centering
\caption{\label{tab:detector} Detector sensitivity comparison for the
PGT bounce GW spectrum. The minimum gap across all frequencies and
parameter choices is $10^{17}$.}
\begin{ruledtabular}
\begin{tabular}{lccc}
\textbf{Detector} & \textbf{Band [Hz]} &
  $\OmGW^{\rm sens}\,h^2$ & \textbf{Min.\ gap} \\
\hline
PTA & $10^{-9}$--$10^{-7}$ & $10^{-10}$ & $10^{42}$ \\
LISA & $10^{-4}$--$10^{-1}$ & $10^{-13}$ & $10^{35}$ \\
DECIGO & $10^{-1}$--$10^{1}$ & $10^{-16}$ & $10^{26}$ \\
LIGO O5 & $10^{1}$--$10^{3}$ & $10^{-10}$ & $10^{20}$ \\
ET & $10^{0}$--$10^{4}$ & $10^{-13}$ & $\mathbf{10^{17}}$ \\
HF concept & $10^{8}$--$10^{10}$ & $10^{-10}$ & $10^{4}$ \\
\end{tabular}
\end{ruledtabular}
\end{table}

The best case is the Einstein Telescope (ET) at $\sim\!10^4$~Hz, with
optimal $\mT \sim 10^7$~GeV giving $\OmGW^{\rm peak}\,h^2 \sim 3\times
10^{-30}$ against a sensitivity of $\sim\!10^{-13}$. The gap is
$10^{17}$---insurmountable.

\subsection{Comparison with other GW sources}

The bounce spectrum (flat + exponential cutoff) differs qualitatively
from inflation (near-flat power law), phase transitions (peaked with
$f^3$ rise and $f^{-1}$ decay), and cosmic strings (flat without cutoff).
However, these shape differences are moot: the amplitude is $\geq 10^{17}$
below any detector threshold. Furthermore, the amplitude scaling
$\OmGW \propto (H_{\rm bounce}/\MPl)^2$ applies equally to LQC,
ekpyrotic, and any other bounce model relying on vacuum amplification.

\subsection{Barrier 12: Vacuum amplification ceiling}

The GW energy density from vacuum amplification at a cosmological bounce
satisfies
\begin{equation}\label{eq:ceiling}
  \OmGW \propto \Bigl(\frac{H_{\rm bounce}}{\MPl}\Bigr)^{\!2}
  \propto \Bigl(\frac{f_b}{f_{\rm Pl}}\Bigr)^{\!4}.
\end{equation}
This fundamental bound ensures that any bounce whose features fall in a
GW detector band ($f < 10^4$~Hz) produces a signal at least $10^{17}$
below sensitivity. The bound is model-independent and applies to all
bounce mechanisms relying on vacuum amplification.

The only escape would be a \textit{non-vacuum} GW source: a pre-bounce
contraction phase that coherently amplifies tensor modes, post-bounce
phase transitions, or topological defect formation. All of these either
require additional physics beyond the bounce or are not bounce-specific.


%% ============================================================
%%  13. BARYOGENESIS AND RELIC PRODUCTION (BRANCH N)
%% ============================================================
\section{Baryogenesis and Relic Production}\label{sec:baryogenesis}

The spin-torsion bounce couples directly to fermion chirality through the
$(J^5)^2$ contact interaction. At the bounce energy scale ($T \sim \MPl$
for ECH, $T \sim \mT$ for PGT), particle production and asymmetry
generation are kinematically open. This motivates testing whether the
bounce can drive baryogenesis, dark matter production, or other relic
mechanisms.

\subsection{Candidate mechanisms screened}

Seven candidate mechanisms were constructed and screened:
\begin{enumerate}[leftmargin=*,label=\Alph*.]
  \item \textit{Axial chemical potential baryogenesis:} The torsion
    $(J^5)^2$ interaction generates an effective chemical potential
    $\mu_5 \sim G\,\xi\,n_5$ for chiral fermions, biasing sphaleron
    processes.
  \item \textit{Modified heavy fermion decay:} Bounce-era curvature
    shifts heavy particle masses by $\delta m \sim \mathcal{O}(\MPl)$,
    modifying CP-violating decay rates.
  \item \textit{Gravitational superheavy dark matter:} Nonadiabatic
    particle production at the bounce creates superheavy relics with
    $m \sim \MPl$.
  \item \textit{Bounce-modified axion abundance:} The bounce modifies the
    axion misalignment angle $\theta_i$.
  \item \textit{Torsion-assisted leptogenesis:} The $(J^5)^2$ interaction
    enhances lepton-number-violating processes.
  \item \textit{PBH formation:} Bounce-scale perturbations collapse to
    primordial black holes.
  \item \textit{Sterile neutrino production:} Spin-density coupling
    populates sterile neutrinos during the bounce.
\end{enumerate}

\subsection{Results}

\textit{All seven candidates are killed.}

The $(J^5)^2$ interaction conserves both baryon number $B$ and lepton
number $L$. It is a chirality-sensitive \textit{scattering} operator, not
a $B$- or $L$-violating operator. Any baryogenesis mechanism must rely on
external $B/L$-violating physics (sphalerons, Majorana masses), with
torsion providing at most a perturbative correction.

At $T \sim \MPl$, all gravitational-strength interactions have comparable
rates: $\Gamma \sim G^2 T^5 \sim H$. The torsion $(J^5)^2$ interaction is
one of approximately $g_* \sim 100$ such channels, contributing
$\sim\!1\%$ of any effect. It cannot dominate any relic abundance.

The strongest candidate (A: axial $\mu_5$ baryogenesis) produces
$\eta_B \sim 10^{-9}$ for maximal initial chirality---tantalizingly close
to observation. But the initial chirality $n_5$ is a \textit{free parameter}
that the bounce does not set (the gravitational chiral anomaly vanishes:
$\tilde{R}R = 0$ on FRW, established in Sec.~\ref{sec:tensor}). The
weakest candidate (D: axion) changes the misalignment angle by
$\Delta\theta \sim 10^{-96}$~rad---96 orders of magnitude below relevance.

\subsection{Barrier 13a: Gravitational democracy}

\textit{At $T \sim \MPl$, all gravitational-strength interactions have
comparable rates. The torsion $(J^5)^2$ interaction contributes
$\sim\!1/g_*$ of the total gravitational effect. It conserves $B$ and $L$,
so it cannot generate asymmetries that generic gravitational interactions
do not. No relic abundance or asymmetry is dominated by, or specific to,
the torsion interaction.}

The bounce is trapped: too high-energy for late-time observables (scale
separation, Branches~A--M) and too generic at its own energy scale for
early-time observables (gravitational democracy, Branch~N).


%% ============================================================
%%  14. HIDDEN-SECTOR VACUUM SELECTION (BRANCH O)
%% ============================================================
\section{Hidden-Sector Vacuum Selection}\label{sec:hidden}

Branch~J established that purely Hamiltonian (reversible) state selection
is killed by Liouville's theorem (Barrier~9). The remaining logical
possibility is \textit{irreversible} mechanisms: thermalization, particle
production, tunneling, dissipation, or phase-transition hysteresis. We
test whether the bounce can trigger such processes in a hidden sector,
indirectly selecting the late-time vacuum energy.

\subsection{Candidate mechanisms screened}

Seven candidates were constructed:
\begin{enumerate}[leftmargin=*,label=\Alph*.]
  \item \textit{First-order phase transition:} Bounce curvature
    ($R \sim \MPl^2$) triggers a hidden-sector transition.
  \item \textit{Metastable vacuum trapping:} Bounce temporarily deforms
    the hidden-sector potential, trapping the field.
  \item \textit{Symmetry restoration/re-breaking:} The bounce restores a
    hidden symmetry that re-breaks with biased branch selection.
  \item \textit{Tunneling enhancement:} Bounce-era potential distortion
    enhances or suppresses quantum tunneling.
  \item \textit{Nonadiabatic production:} Bounce populates a pNGB hidden
    sector via particle production.
  \item \textit{Dissipative reheating:} Post-bounce thermalization
    selects a thermal branch.
  \item \textit{Domain-wall percolation:} Bounce-imposed asymmetry biases
    multi-vacuum domain percolation.
\end{enumerate}

\subsection{Results}

\textit{All seven candidates are killed.}

The results partition into three structural outcomes:

\textit{(i) Irreversibility is real but not bounce-specific.} Five
candidates (B, C, E, F, G) involve genuine irreversibility, but the same
processes occur in any hot cosmology. The bounce provides energy, but so
does inflationary reheating.

\textit{(ii) Bounce-specific candidates fail on naturalness.} Two
candidates (A, D) have marginal bounce-specificity: only the
Planck-scale curvature of the ECH bounce can trigger the transition. But
the resulting vacuum energy $\rho_{\rm DE}$ is set by hidden-sector
potential parameters, not by the bounce. The bounce determines
\textit{whether} a transition occurs, not \textit{what} vacuum energy
results.

\textit{(iii) The $10^{122}$ hierarchy is untouched.} No mechanism
connects $\MPl^4$ to $(2.3\;\text{meV})^4$ through bounce dynamics. The
hierarchy must be bridged by hidden-sector parameters that the bounce
does not set.

\subsection{Barrier 13b: Bounce--vacuum energy decoupling}

\textit{The bounce operates at scale $M$ ($= \MPl$ or $\mT$). The
late-time vacuum energy is at scale $\Lambda_{\rm DE}^4$. Any formula
$\rho_{\rm DE} = f(M, g_i)$ requires a suppression factor
$\sim\!(\Lambda_{\rm DE}/M)^p$ that cannot be generated by bounce
dynamics. Adding irreversibility determines whether a transition occurs
but not the vacuum energy scale. The trigger and the outcome are
structurally decoupled.}

\subsection{Completeness of the state-change closure}

Barriers~9 and~13 together close \textit{all} state-change mechanisms:
\begin{itemize}[leftmargin=*]
  \item Reversible (Hamiltonian) mechanisms: killed by Liouville
    (Barrier~9).
  \item Irreversible, bounce-specific mechanisms: killed by decoupling
    (Barrier~13).
  \item Irreversible, non-bounce-specific mechanisms: not relevant to the
    bounce program.
\end{itemize}
There is no remaining category. The logical space of bounce$\to$DE
state-change connections is exhausted.


%% ============================================================
%%  15. PGT TORSION RELIC COSMOLOGY (BRANCH P)
%% ============================================================
\section{PGT Torsion Relic Cosmology}\label{sec:relic}

The last surviving positive channel was torsion relic cosmology: if the
PGT Sector~II pseudoscalar torsion mode $\phi(t) \equiv S_0(t)$ carries
$\mathcal{O}(1)$ of the bounce energy, its post-bounce decay yields a
massive relic constrained by BBN, CMB spectral distortions, and dark
radiation bounds. The gating question is whether $\phi$ is actually
populated at the bounce.

\subsection{FRW parity protection}

The Sector~II pseudoscalar mode $S_0(t)$ is kinematically allowed on FRW:
isotropy permits $S_\mu = (S_0(t), 0, 0, 0)$. However, $S_0$ is
parity-\textit{odd} (a pseudoscalar under spatial inversion). FRW
spacetime is parity-invariant, and the PGT Sector~II action is quadratic
in $S_\mu$, possessing an exact $\mathbb{Z}_2$ symmetry $\phi \to -\phi$.
The field equation on FRW is
\begin{equation}\label{eq:torsion_kg}
  \ddot{\phi} + 3H\dot{\phi} + \mT^2 \phi = 0\,,
\end{equation}
with \textit{no source term}: parity-even matter cannot source a
parity-odd field. The origin $\phi = 0$ is a stable fixed point
($\mT^2 > 0$), and the bounce---being a parity-even gravitational
event---cannot dislodge it.

\subsection{Population mechanisms tested}

Five population scenarios were assessed:
\begin{enumerate}[leftmargin=*]
  \item \textit{Self-consistent bounce:} The modified Friedmann equation
    works with $\phi = 0$ (the contact-term $(J^5)^2$ alone drives the
    bounce). The propagating mode is not required.
  \item \textit{Quantum vacuum fluctuations:} Produce particle pairs
    (occupation number $\langle n_k \rangle$), not a homogeneous
    condensate $\langle\phi\rangle \neq 0$.
  \item \textit{Chiral matter sourcing:} A chiral asymmetry $n_5 \neq 0$
    in the matter sector could source $\phi$ through the axial coupling,
    but $n_5$ is a free parameter (the gravitational chiral anomaly
    $\tilde{R}R = 0$ on FRW; see Sec.~\ref{sec:tensor}) and the sourced
    amplitude is $\phi \sim G\,n_5 \ll \MPl$.
  \item \textit{Parametric resonance:} The bounce is $\sim\!1$ oscillation
    period long; no time for resonant buildup.
  \item \textit{Unmotivated initial condition:} $\phi(0) \sim \MPl$ as a
    free parameter. Not a prediction---equivalent to assuming the
    conclusion.
\end{enumerate}
No mechanism populates the pseudoscalar mode at $\mathcal{O}(\MPl)$ or
$\mathcal{O}(\mT)$.

\subsection{Literature context}

The PGT cosmology literature (Shie, Nester, Yo 2008; Chen \etal\ 2009)
uses the parity-\textit{even} trace vector $T_0(t)$, which is naturally
nonzero on FRW. This is a fundamentally different mode (Sector~I/III),
which has ghost issues at the target mass scale. The specific question
of pseudoscalar population at a PGT bounce had not been addressed; our
answer is negative.

\subsection{Barrier 14: Parity protection of the torsion vacuum}

\textit{The PGT Sector~II pseudoscalar torsion mode has an exact
$\mathbb{Z}_2$ parity symmetry $\phi \to -\phi$ that protects $\phi = 0$
as the unique parity-invariant vacuum on FRW. No parity-even dynamics
(including the bounce) can excite the mode. The torsion relic abundance
is zero unless parity is explicitly broken by the initial conditions.}

This closes the last positive observable channel of the minimal PGT
bounce program.


%% ============================================================
%%  16. COMPLETE BARRIER CATALOG
%% ============================================================
\section{Complete Barrier Catalog}\label{sec:barriers}

Fourteen structural barriers have been identified across the full research
program (Branches~A--P). They span five distinct failure modes:
energy-budget failure (Barriers~1--7), signature failure
(Barriers~8, 10--12), reversible state-selection failure (Barrier~9),
irreversible state-selection / relic failure (Barrier~13), and relic
population failure (Barrier~14).

\begin{table*}[t]
\centering
\caption{\label{tab:barriers} Complete catalog of fourteen structural
barriers identified in the spin-torsion bounce and dark energy program.}
\begin{ruledtabular}
\begin{tabular}{clll}
\textbf{\#} & \textbf{Barrier} & \textbf{Branch} & \textbf{Description} \\
\hline
1 & Mass-coupling lock & Found.\ A (PGT) & $m_B$ and $\geff$ decrease together; $m_B/\geff \sim \MPl^2$ fixed \\
2 & Topological-shift duality & Found.\ B & Nieh-Yan non-topological in MAG but geometric content blocked \\
3 & Scalar-tensor universality & Found.\ C & Environmental mass reduces to standard scalar-tensor on FRW \\
4 & Planck suppression & Found.\ D & Connection coupling gives 1~$\partial\phi$/vertex; disformal requires 2 \\
5 & Scale separation & Found.\ E & $k_{\rm CMB}/k_b \sim 10^{-28}$; global integrals insensitive to bounce \\
6 & Attractor-sensitivity dilemma & Found.\ F & Attractors erase IC; sensitivity requires fine-tuning \\
7 & Parameter immunity & Found.\ G & Cyclic incompatibility; continuous vacuum solutions \\
8 & Parity-even effective interaction & Branch H & $(J^5)^2$ is parity-even; FRW kills spatial P-odd backgrounds \\
9 & Hamiltonian phase-space conservation & Branch J & Liouville prevents phase-space contraction at bounce \\
10 & UV$\to$IR specificity dilemma & Branch L$_1$ & Observable $\Rightarrow$ generic; specific $\Rightarrow$ non-minimal \\
11 & Decoupling universality & Branch L$_2$ & Light PGT gauge fields decouple universally \\
12 & Vacuum amplification ceiling & Branch M & $\OmGW \propto (H/\MPl)^2$; minimum $10^{17}$ gap to detectors \\
13 & Gravitational democracy / & Branches N, O & At $T\sim\MPl$, torsion is 1 of $\sim\!100$ channels; \\
   & bounce--vacuum decoupling & & bounce trigger decoupled from vacuum energy outcome \\
14 & Parity protection of torsion vacuum & Branch P & $\mathbb{Z}_2$: $\phi\to-\phi$ protects $\phi=0$; no population mechanism \\
\end{tabular}
\end{ruledtabular}
\end{table*}


%% ============================================================
%%  14. DISCUSSION
%% ============================================================
\section{Discussion}\label{sec:discussion}

\subsection{The generic-specific dichotomy}

The central finding of the bounce characterization is a sharp dichotomy:
\textit{what is detectable is generic to any bounce; what is specific to
torsion is undetectable.}

Observable quantities---the flat tensor spectrum ($n_T \approx 0$), the
unit transfer function ($T(k) = 1$), the GW spectral shape (flat +
exponential cutoff)---are all shared by any time-symmetric radiation
bounce. They are properties of the bounce \textit{geometry}, not the
bounce \textit{mechanism}.

Torsion-specific quantities---the critical density $\rhocrit \approx
0.21\,\MPl^4$, the exact bounce profile, the $(J^5)^2$ corrections, the
chirality $\Delta\chi = 0$---are all either unobservable (confined to
GHz scales) or null (parity-even interaction, Pontryagin density vanishes
on FRW).

\subsection{Comparison with other bouncing cosmologies}

\subsubsection{Loop quantum cosmology (LQC)}

LQC also produces a modified Friedmann equation of the
form~\eqref{eq:modfriedmann} with $\rhocrit \approx 0.41\,\rhoPl$
(slightly different numerical coefficient). The tensor spectrum has the
same functional form: flat, with $P_T \sim 10^{-64}$ and $n_T \approx
0$. The GW amplitude-frequency tradeoff~\eqref{eq:ceiling} applies
identically. The spin-torsion bounce and LQC bounce are observationally
indistinguishable at any accessible frequency.

\subsubsection{Ekpyrotic bounce}

Ekpyrotic scenarios produce a strongly blue-tilted tensor spectrum
($n_T \sim 2$), in contrast to the flat spectrum of the spin-torsion and
LQC bounces. However, the ekpyrotic bounce amplitude is similarly
suppressed. In principle, $n_T$ distinguishes ecpyrotic from
radiation-symmetric bounces---but only if the amplitude is detectable,
which it is not.

\subsubsection{Matter bounce}

The matter bounce scenario~\cite{Quintin2015,Cai2014} also produces
$n_T \approx 0$ (scale-invariant tensors), making it degenerate with
the spin-torsion bounce. The matter bounce requires a specific
pre-bounce contraction phase (dust-like), which is not part of the
minimal spin-torsion model. Both models face the same amplitude
suppression.

\subsection{The complementary trap}

The spin-torsion bounce is trapped between three complementary failures:
\begin{itemize}[leftmargin=*]
  \item \textit{Too high-energy} for late-time observables (scale
    separation, Planck suppression---Branches A--M).
  \item \textit{Too generic} at its own energy scale (gravitational
    democracy---Branch~N). At $T \sim \MPl$, torsion is one
    gravitational-strength channel among $\sim\!100$.
  \item \textit{Too decoupled} from late-time vacuum energy
    (bounce--vacuum decoupling---Branch~O). The bounce trigger and the
    vacuum energy outcome are structurally independent.
\end{itemize}

\subsection{What would change the assessment}

The assessment would change if:
\begin{enumerate}[leftmargin=*]
  \item A \textit{non-vacuum} GW source is identified at the bounce
    (particle production, phase transition) that does not suffer from
    the vacuum amplification ceiling. This would require a detailed model
    of the bounce-to-inflation transition.
  \item A \textit{pre-bounce} contraction phase with $w \neq 1/3$ is
    specified, allowing horizon exit of tensor modes at CMB scales during
    contraction. This is additional physics beyond the spin-torsion bounce.
  \item A mechanism breaking the mass-coupling lock is found (independent
    mass-generating mechanism), enabling observable torsion-specific
    signatures.
  \item Non-perturbative or multi-loop effects generate $\BI$-dependent
    vacuum structure not accessible at the tested orders.
  \item A mechanism is found that connects the bounce energy scale to the
    vacuum energy scale without fine-tuning, evading both Liouville
    (Barrier~9) and bounce--vacuum decoupling (Barrier~13).
\end{enumerate}
None of these outcomes is assumed or expected within the minimal and
near-minimal program assessed here.

\subsection{The program as a publishable negative result}

Negative results in theoretical physics are undervalued but scientifically
important. The combined program (four minimal DE routes, seven extended
foundations, five bounce observable channels, seven baryogenesis/relic
candidates, seven hidden-sector vacuum candidates) establishes:
\begin{itemize}[leftmargin=*]
  \item A complete classification of failure modes (energy-budget,
    signature, reversible state-selection, irreversible state-selection /
    relic production).
  \item Fourteen structural barriers spanning five failure modes.
  \item Six structural lessons and five decision requirements (DR1--DR5)
    constraining future proposals.
  \item A reproducible methodology for disciplined model testing.
  \item Exhaustive coverage of the logical space: direct derivation
    (A--G), observable signatures (H, K, L--M), state selection (J, O),
    UV$\to$IR bridge (L), relic production (N), and GW spectrum (M).
\end{itemize}
The program's value lies not in any single null result but in the
\textit{completeness} of the closure: every natural direction has been
tested, and the structural reasons for failure have been identified and
cataloged.


%% ============================================================
%%  15. CONCLUSION
%% ============================================================
\section{Conclusion}\label{sec:conclusions}

This paper presents a comprehensive theoretical assessment of the
Einstein-Cartan-Holst framework as both a dark-energy phenomenology and a
bouncing cosmology. The contribution is threefold.

\textit{First}, the phenomenological dark-energy ansatz rooted in ECH
gravity remains observationally consistent. It reduces fine-tuning from
$10^{120}$ to $\sim\!10^5$, yields MCMC fits partially reducing the $H_0$
and $\sigma_8/S_8$ tensions, and has a parity-odd structure qualitatively
consistent with observed birefringence. The equation of state $w=-1$ is
assumed, not derived.

\textit{Second}, the systematic closure of four minimal and seven extended
derivation routes establishes that neither the minimal ECH model nor its
most natural extensions can derive late-time dark energy from geometry.
Seven structural barriers and six structural lessons map the failure
landscape. The mass-coupling lock (Barrier~1)---in which a light torsion
mass automatically decouples the field from matter---is a previously
unrecognized obstruction that applies to any geometric theory where mass
and coupling share a common kinetic-normalization origin.

\textit{Third}, the comprehensive characterization of the spin-torsion
bounce across all observable channels establishes its observational
inertness. The tensor spectrum is $P_T \approx 2\times 10^{-64}$
(undetectable). The chirality is $\Delta\chi = 0$ exactly (no parity
violation). The scalar transfer function is $T(k) = 1$ (no bounce
imprint). State selection is prevented by Liouville's theorem. The PGT
gravitational wave spectrum has a minimum $10^{17}$ gap to any detector.
Baryogenesis and relic production via the $(J^5)^2$ interaction are
killed by gravitational democracy at $T \sim \MPl$. Hidden-sector vacuum
selection via irreversible transitions is killed by bounce--vacuum energy
decoupling: the bounce determines \textit{whether} a transition occurs,
not \textit{what} vacuum energy results. The last positive channel---PGT
torsion relic cosmology---is closed by $\mathbb{Z}_2$ parity protection:
the pseudoscalar mode cannot be populated at the bounce without explicit
parity violation in the initial conditions. Seven additional barriers
(8--14) are cataloged.

The spin-torsion bounce is theoretically consistent but observationally
inert. Fourteen structural barriers, spanning ten research branches and
five failure modes, prevent any nontrivial connection between the
Planck-scale bounce and low-energy cosmological observations. The logical
space of bounce$\to$observable connections is exhausted: direct derivation,
observable signatures, state selection (both reversible and irreversible),
relic production, torsion relic population, and UV$\to$IR bridges have all
been tested and closed. Future directions require physics beyond minimal
Einstein-Cartan gravity and quadratic Poincar\'e gauge theory.


%% ============================================================
%%  ACKNOWLEDGMENTS
%% ============================================================
\begin{acknowledgments}
The author thanks the anonymous referees for constructive feedback.
Computational resources for MCMC analyses were provided by RunPod.
This work made use of Cobaya~\cite{Cobaya2021} and CAMB.
\end{acknowledgments}


%% ============================================================
%%  APPENDICES
%% ============================================================
\appendix

%% ============================================================
%%  A. NOTATION AND CONVENTIONS
%% ============================================================
\section{Notation and Conventions}\label{app:notation}

Natural units $c=\hbar=1$ throughout, with $8\pi G = \MPl^{-2}$.
Metric signature $(-,+,+,+)$. Reduced Planck mass
$\MPl = (8\pi G)^{-1/2} = 2.435 \times 10^{18}$~GeV.
$\gamma^5 = i\gamma^0\gamma^1\gamma^2\gamma^3$ with
$(\gamma^5)^2 = 1$. The vacuum energy density is
$\rho_\Lambda = \Leff\,\MPl^2/(8\pi)$; when we write
$\rho_\Lambda \approx (2.3\;\text{meV})^4$ this refers to the energy
density. Greek indices $\mu,\nu,\ldots$ denote spacetime coordinates;
Latin indices $I,J,\ldots$ denote tangent-space (Lorentz) indices.
The Levi-Civita symbol $\varepsilon^{IJKL}$ is totally antisymmetric
with $\varepsilon^{0123}=+1$.

Key symbols:
\begin{itemize}[leftmargin=*]
\item $\BI$---Barbero-Immirzi parameter ($0.274 \pm 0.020$)
\item $K^{IJ}_{\;\;\;\mu}$---contorsion tensor
\item $T^I_{\;\;JK}$---torsion tensor
\item $\alpha/M$---parity-odd coefficient (mass dimension $-1$)
\item $\Dinf$---inflationary dilution factor ($\sim e^{-3N_{\rm tot}}$)
\item $\Xi\equiv[(\alpha/M)\,\MPl]\,\Dinf$---inflationary suppression
  factor; $\rho_\Lambda = \Xi\,\MPl^4$
\item $\Geff = (3\kappa^2/16)\,\BI^2/(\BI^2+1)$---effective four-fermion
  coupling
\item $\Gsp \propto (\BI^2-1)/(\BI^2+1)$---scalar/pseudoscalar channel
  coupling
\item $\mT$---PGT torsion mass
\item $\geff$---canonically normalized matter coupling
\item $\rhocrit$---bounce critical density
\item $k_b$---bounce characteristic wavenumber
\item $f_b$---bounce characteristic frequency
\end{itemize}

Abbreviations: LQG (Loop Quantum Gravity), EC (Einstein-Cartan),
ECH (Einstein-Cartan-Holst), PGT (Poincar\'e Gauge Theory),
ALP (axion-like particle), NJL (Nambu-Jona-Lasinio), MAG
(metric-affine gravity), MCMC (Markov Chain Monte Carlo), ET (Einstein
Telescope), HF (high-frequency).


%% ============================================================
%%  B. PARAMETER SUMMARY
%% ============================================================
\section{Parameter Summary}\label{app:params}

\begin{table*}[t]
\caption{\label{tab:params} Complete parameter classification. Status
categories: \textit{Fixed} (external input), \textit{Ansatz}
(phenomenological scaling), \textit{Assumed} (not derived),
\textit{Fit} (MCMC), \textit{Derived} (from other parameters).}
\begin{ruledtabular}
\begin{tabular}{llcll}
\textbf{Parameter} & \textbf{Description} & \textbf{Value}
  & \textbf{Status} & \textbf{Notes} \\
\hline
\multicolumn{5}{l}{\textit{Framework parameters}} \\
$\BI$ & Barbero-Immirzi & $0.274 \pm 0.020$ & Fixed & LQG area spectrum \\
$\alpha/M$ & Parity-odd coeff. & $\sim 10^{-33}\;\text{eV}^{-1}$
  & Fit & One-loop motivated \\
$N_{\rm tot}$ & Total $e$-folds & $\sim 92$ & Assumed
  & Fitted to $\rho_\Lambda$ \\
$\Dinf$ & Inflation dilution & $\sim e^{-3N_{\rm tot}}$
  & Derived & From $N_{\rm tot}$ \\
$\Xi$ & Infl.\ suppression & $\sim 10^{-123}$
  & Ansatz & $=[(\alpha/M)\MPl]\Dinf$ \\
$w$ & Equation of state & $-1$ & Assumed & Not derived \\
\hline
\multicolumn{5}{l}{\textit{Standard cosmological parameters (MCMC)}} \\
$H_0$ & Hubble constant & $69.2\pm 0.8$\,\textsuperscript{a}
  & Fit & km/s/Mpc \\
$\sigma_8$ & Clustering ampl. & $0.785\pm 0.016$\,\textsuperscript{a}
  & Fit & \\
$S_8$ & $\sigma_8(\Omega_m/0.3)^{0.5}$ & $0.80\pm 0.02$
  & Derived & \\
$\Omega_m$ & Matter density & $0.310 \pm 0.008$ & Fit & \\
$\Omega_b h^2$ & Baryon density & $0.02237 \pm 0.00015$ & Fit & \\
$n_s$ & Spectral index & $0.9649 \pm 0.0042$ & Fit & \\
$\tau$ & Optical depth & $0.054 \pm 0.007$ & Fit & \\
\hline
\multicolumn{5}{l}{\textit{Extended/bounce parameters}} \\
$\Delta\Neff$ & Extra species & $\sim 0$\,\textsuperscript{a}
  & Fit & Consistent with 0 \\
$\rhocrit$ & Bounce density & $0.21\,\MPl^4$ & Derived & Minimal ECH \\
$\mT$ & PGT torsion mass & Free & Model param. & Sets PGT bounce scale \\
$|t_3|$ & PGT coupling & Free & Model param. & $\mT = \MPl/(2\sqrt{|t_3|})$ \\
\end{tabular}
\end{ruledtabular}
\textsuperscript{a}Full tension dataset (includes SH0ES $H_0$ prior).
Verification finds $H_0 = 67.68 \pm 1.06$, $\Delta\Neff = -0.020
\pm 0.169$.
\end{table*}


%% ============================================================
%%  C. CLAIMS CLASSIFICATION
%% ============================================================
\section{Claims Classification}\label{app:claims}

\begin{table*}[t]
\caption{\label{tab:claims} Expanded claims classification. Categories:
Derived (structural results), Phenomenological (ansatz/fit/assumption),
Negative (clean null results from the testing program), and
Retired (superseded by closures).}
\begin{ruledtabular}
\begin{tabular}{llll}
\textbf{Claim} & \textbf{Status} & \textbf{Basis} & \textbf{Notes} \\
\hline
\multicolumn{4}{l}{\textit{Derived (structural results)}} \\
ECH action structure & Derived & Standard EC + Holst & Background \\
Torsion is algebraic (minimal ECH) & Derived & EC field equations & Structural \\
Perfect-square four-fermion & Derived & Torsion elimination & Exact result \\
$\Gsp < 0$ at $\BI < 1$ & Derived & Fierz rearrangement & Track B closure \\
One-loop $\slashed{D}$ is $\BI$-free & Derived & Torsion elimination & G v1 closure \\
Dyn.\ Immirzi $=$ ALP & Derived (lit.) & Calcagni-Mercuri 2009 & T1 closure \\
No photon coupling in min.\ ECH & Derived & Action structure & S1 closure \\
Nieh-Yan is topological in 4D & Derived (lit.) & Differential geometry & T1 closure \\
$(J^5)^2$ is parity-even & Derived & Algebra & Branch H closure \\
$T(k) = 1$ (time-reversal) & Derived & Analytic + numerical & Branch K \\
$\Delta\chi = 0$ on FRW & Derived & FRW isotropy & Branch H closure \\
Liouville prevents state selection & Derived & Hamiltonian mechanics & Branch J closure \\
$\OmGW \propto (\mT/\MPl)^2$ & Derived & Bogoliubov calculation & Branch M \\
\hline
\multicolumn{4}{l}{\textit{Phenomenological (ansatz, fit, or assumption)}} \\
$\rho_\Lambda = \Xi\,\MPl^4$ & Ansatz & Scaling argument & Not a derivation \\
$w = -1$ at late times & Assumed & Observational input & Not derived \\
$H_0$, $\sigma_8$, $S_8$ values & Fit & MCMC & Standard analysis \\
Fine-tuning $10^{120} \to 10^5$ & Parametric & Sensitivity scan & Valid \\
\hline
\multicolumn{4}{l}{\textit{Negative results (clean null)}} \\
$P_T \approx 2\times 10^{-64}$ & Null (amplitude) & Numerical & Undetectable \\
$n_T \approx 0$ (generic) & Null (shape) & Numerical + analytic & Any rad.\ bounce \\
$r \sim 10^{-55}$ & Null (ratio) & Derived & Undetectable \\
$f_{\rm NL} \sim 10^{-56}$ & Null (NG) & Order-of-magnitude & Undetectable \\
GW gap $\geq 10^{17}$ & Null (GW) & Bogoliubov + detector & Branch M \\
No Horndeski constraint & Null (compat.) & Scale separation & Branch I \\
\hline
\multicolumn{4}{l}{\textit{Retired (closed by testing program)}} \\
Birefringence as prediction & Retired & S1 closure & No photon coupling \\
$\BI$ drives late-time physics & Retired & B, G, T1 closures & UV-scale only \\
Condensate mechanism & Retired & Track B closure & $\Gsp < 0$ \\
PGT torsion as DE mechanism & Retired & Foundation A & Mass-coupling lock \\
Bounce tensor as observable & Retired & Branch H & $P_T \sim 10^{-64}$ \\
Bounce chirality as observable & Retired & Branch H & $\Delta\chi = 0$ \\
Bounce state selection & Retired & Branch J & Liouville theorem \\
PGT GW as observable & Retired & Branch M & Vac.\ ampl.\ ceiling \\
\end{tabular}
\end{ruledtabular}
\end{table*}


%% ============================================================
%%  D. BARDEEN EQUATION THROUGH THE BOUNCE
%% ============================================================
\section{Bardeen Equation Through the Bounce}\label{app:bardeen}

We sketch the derivation of the Bardeen equation on the spin-torsion
bounce background and the proof of $T(k) = 1$.

\subsection{Setup}

Consider scalar perturbations of a radiation fluid ($w = 1/3$,
$c_s^2 = 1/3$) on the modified Friedmann
background~\eqref{eq:modfriedmann}. In the longitudinal gauge
($B = E = 0$), the perturbed metric is
\begin{equation}
  ds^2 = -(1+2\Phi)\,dt^2 + a^2(1-2\Psi)\,\delta_{ij}\,dx^i\,dx^j,
\end{equation}
with $\Phi = \Psi$ for a perfect fluid (no anisotropic stress).

\subsection{Derivation of the Bardeen equation}

From the linearized Einstein equations on the modified Friedmann
background, one obtains~\cite{MukhanovFeldmanBrandenberger1992,
PeterPintoNeto2008}:
\begin{equation}\label{eq:bardeen_app}
  \ddot\Phi + 5H\dot\Phi +
  \Bigl[\frac{k^2}{3a^2} + 2\dot H + 4H^2\Bigr]\Phi = 0.
\end{equation}
The key feature for bouncing cosmologies: all coefficients are
\textit{regular} at $H = 0$. Specifically:
\begin{itemize}[leftmargin=*]
  \item The friction term $5H\dot\Phi$ vanishes at $H = 0$ (the bounce
    instant).
  \item The effective mass term $2\dot H + 4H^2$ remains finite:
    $\dot H(0) \approx 3.52\,\MPl^2$ and $H(0) = 0$.
  \item The gradient term $k^2/(3a^2)$ is regular and positive.
\end{itemize}
This regularity is in contrast to the Mukhanov-Sasaki variable
$v = z\mathcal{R}$ where $z = a\dot\phi/H \to \infty$ at the bounce.

\subsection{Time-reversal symmetry and $T(k) = 1$}

The bounce background $a(t)$ is an even function of $t$:
$a(-t) = a(t)$. Therefore $H(-t) = -H(t)$ and $\dot H(-t) = \dot H(t)$.
Under the substitution $t \to -t$, Eq.~\eqref{eq:bardeen_app} maps to
\begin{equation}
  \ddot\Phi - 5H\dot\Phi +
  \Bigl[\frac{k^2}{3a^2} + 2\dot H + 4H^2\Bigr]\Phi = 0,
\end{equation}
where the only change is the sign of the friction term. This means that
if $\Phi_{\rm grow}(t)$ is a growing solution in the contracting phase
($t < 0$, $H < 0$, friction term is anti-damping), then
$\Phi_{\rm grow}(-t)$ is a \textit{decaying} solution in the expanding
phase ($t > 0$).

The constant mode $\Phi_{\rm const}$ (the physically relevant mode for
the power spectrum) is symmetric under $t \to -t$ and maps to itself
with unit amplitude. Therefore:
\begin{equation}
  T(k) \equiv \frac{\Phi_{\rm const}^{\rm out}}
  {\Phi_{\rm const}^{\rm in}} = 1
  \quad \text{for all } k \ll k_b.
\end{equation}


%% ============================================================
%%  E. PGT GHOST-FREE CONDITIONS
%% ============================================================
\section{PGT Ghost-Free Conditions}\label{app:ghost}

\subsection{Quadratic PGT action}

The most general quadratic PGT action with torsion-squared and
curvature-squared terms has ten free parameters (three torsion-squared
$a_1, a_2, a_3$ plus six curvature-squared $b_i$ plus
$\Lambda_0$)~\cite{HayashiShirafuji1980,Neville1980,
SezginNieuwenhuizen1980}.

\subsection{Single-mode ghost-free models}

Sezgin and van Nieuwenhuizen~\cite{SezginNieuwenhuizen1980} identified
the conditions for ghost-freedom (absence of negative-norm states) in
the linearized theory around flat spacetime. Karananas~\cite{Karananas2015}
confirmed and extended this analysis. The cleanest single-mode models are:

\begin{itemize}[leftmargin=*]
  \item \textit{Spin-$0^-$ (axial pseudoscalar):} Ghost-free for all
    $|t_3| > 0$. Mass $m_B = \MPl/(4\sqrt{\pi|t_3|})$. This is the mode
    assessed in Foundation~A.
  \item \textit{Spin-$0^+$ (scalar):} Ghost-free only in a restricted
    parameter range. Develops a ghost at the mass scale needed for
    observable bounce features. Rejected.
  \item \textit{Spin-$2^+$ (tensor):} Ghost-free but massive spin-2
    fields face the van Dam-Veltman-Zakharov discontinuity and require
    careful Vainshtein screening.
\end{itemize}

\subsection{Sector II analysis}

For the spin-$0^-$ mode in Sector~II, the propagator is
\begin{equation}
  D(p^2) = \frac{1}{p^2 - m_B^2},
\end{equation}
with no ghost pole (residue is positive). The canonically normalized
field $\phi_{\rm can} = \sqrt{Z}\,\phi$ with $Z \propto |t_3|$ gives
\begin{equation}
  m_B = \frac{\MPl}{4\sqrt{\pi|t_3|}}, \qquad
  \geff = \frac{g}{\sqrt{Z}} \sim \frac{1}{\MPl\sqrt{|t_3|}},
\end{equation}
confirming the mass-coupling lock~\eqref{eq:lock}.

\subsection{Symmetry analysis}

Five candidate symmetries for protecting $m_B$ were investigated:
\begin{enumerate}[leftmargin=*]
  \item \textit{Gauge symmetry:} The $0^-$ mode is not a gauge field;
    no gauge invariance protects its mass.
  \item \textit{Shift symmetry:} The mass term $m_B^2\phi^2$ explicitly
    breaks any shift symmetry $\phi \to \phi + c$.
  \item \textit{Chiral symmetry:} The axial current coupling preserves
    chiral symmetry, but this does not protect a scalar mass.
  \item \textit{Topological protection:} No topological mechanism
    applies in 4D for a massive pseudoscalar.
  \item \textit{Supersymmetry:} Would pair the $0^-$ with a fermion,
    but SUSY is not part of PGT and its breaking scale is unknown.
\end{enumerate}
The 't~Hooft naturalness criterion fails: setting $m_B = 0$ does not
restore any symmetry of the PGT action. Graviton loops generate
\begin{equation}
  \delta m_B^2 \sim \frac{\MPl^2}{16\pi^2},
\end{equation}
requiring fine-tuning of order $10^{-57}$ to maintain $m_B \sim
\text{meV}$.


%% ============================================================
%%  REFERENCES
%% ============================================================
\begin{thebibliography}{99}

\bibitem{Weinberg1989}
S.~Weinberg,
``The cosmological constant problem,''
Rev.\ Mod.\ Phys.\ \textbf{61}, 1 (1989).

\bibitem{Planck2018params}
Planck Collaboration (N.~Aghanim \etal),
``Planck 2018 results.\ VI.\ Cosmological parameters,''
Astron.\ Astrophys.\ \textbf{641}, A6 (2020).

\bibitem{Riess2022}
A.~G.~Riess \etal,
``A comprehensive measurement of the local value of the Hubble constant
with 1 km/s/Mpc uncertainty from the Hubble Space Telescope and the
SH0ES team,''
Astrophys.\ J.\ Lett.\ \textbf{934}, L7 (2022).

\bibitem{KiDS2021}
KiDS Collaboration (M.~Asgari \etal),
``KiDS-1000 cosmology: Cosmic shear constraints on the amplitude of
matter fluctuations,''
Astron.\ Astrophys.\ \textbf{645}, A104 (2021).

\bibitem{DES2024}
DES Collaboration (T.~M.~C.~Abbott \etal),
``Dark Energy Survey Year 3 results: Cosmological constraints from
galaxy clustering and weak lensing,''
Phys.\ Rev.\ D~\textbf{105}, 023520 (2022).

\bibitem{DESI2024}
DESI Collaboration,
``DESI 2024 VI: Cosmological constraints from the measurements of
baryon acoustic oscillations,''
arXiv:2404.03002 (2024).

\bibitem{DESI2025DR2}
DESI Collaboration,
``DESI DR2 results I: Baryon acoustic oscillations and cosmological
constraints,''
arXiv:2503.14738 (2025).

\bibitem{Ashtekar2011}
A.~Ashtekar and P.~Singh,
``Loop quantum cosmology: A status report,''
Class.\ Quant.\ Grav.\ \textbf{28}, 213001 (2011).

\bibitem{Freidel2005}
L.~Freidel, D.~Minic, and T.~Takeuchi,
``Quantum gravity, torsion, parity violation and all that,''
Phys.\ Rev.\ D~\textbf{72}, 104002 (2005).

\bibitem{Mercuri2006}
R.~Mercuri,
``Fermion interactions, cosmological constant and torsion,''
Phys.\ Rev.\ D~\textbf{73}, 084016 (2006).

\bibitem{Mercuri2009}
R.~Mercuri,
``Peccei-Quinn mechanism in gravity and the nature of the
Barbero-Immirzi parameter,''
Phys.\ Rev.\ Lett.\ \textbf{103}, 081302 (2009).

\bibitem{Poplawski2012}
N.~J.~Pop\l{}awski,
``Cosmological constant from quarks and torsion,''
Phys.\ Rev.\ D~\textbf{85}, 107502 (2012).

\bibitem{ShapiroTeixeira2014}
I.~L.~Shapiro and P.~M.~Teixeira,
``Quantum Einstein-Cartan theory with the Holst term,''
Class.\ Quant.\ Grav.\ \textbf{31}, 185002 (2014).

\bibitem{Shapiro2002}
I.~L.~Shapiro,
``Physical aspects of the space-time torsion,''
Phys.\ Rept.\ \textbf{357}, 113 (2002).

\bibitem{Vassilevich2003}
D.~V.~Vassilevich,
``Heat kernel expansion: User's manual,''
Phys.\ Rept.\ \textbf{388}, 279 (2003).

\bibitem{CalcagniMercuri2009}
G.~Calcagni and R.~Mercuri,
``Barbero-Immirzi field in canonical formalism of pure gravity,''
Phys.\ Rev.\ D~\textbf{79}, 104003 (2009).

\bibitem{TaverasYunes2008}
V.~Taveras and N.~Yunes,
``Barbero-Immirzi parameter as a scalar field: K-inflation from loop
quantum gravity?''
Phys.\ Rev.\ D~\textbf{78}, 064070 (2008).

\bibitem{Eskilt2022}
J.~R.~Eskilt and E.~Komatsu,
``Improved constraints on cosmic birefringence from the WMAP and Planck
cosmic microwave background polarization data,''
Phys.\ Rev.\ D~\textbf{106}, 063503 (2022).

\bibitem{DiegoPalazuelos2025}
P.~Diego-Palazuelos and E.~Komatsu,
``Cosmic birefringence from ACT DR6,''
(2025).

\bibitem{Golden2026v1}
H.~Golden,
``Geometric Dark Energy from Spin-Torsion Cosmology:
Phenomenological Constraints and Correlated Signatures,''
Preprint HUBIFY-2026-001, v1.6 (2026).

\bibitem{Cobaya2021}
J.~Torrado and A.~Lewis,
``Cobaya: Code for Bayesian analysis of hierarchical physical models,''
JCAP \textbf{05}, 057 (2021).

\bibitem{ECTorsionDESI2025}
J.~Liu, R.~Niu, and R.-G.~Cai,
``Torsion cosmology with DESI DR2,''
arXiv:2503.02599 (2025).

\bibitem{SezginNieuwenhuizen1980}
E.~Sezgin and P.~van Nieuwenhuizen,
``New ghost-free gravity Lagrangians with propagating torsion,''
Phys.\ Rev.\ D~\textbf{21}, 3269 (1980).

\bibitem{Karananas2015}
G.~K.~Karananas,
``The particle spectrum of parity-violating Poincar\'e gravitational
theory,''
Class.\ Quant.\ Grav.\ \textbf{32}, 055012 (2015).

\bibitem{HayashiShirafuji1980}
K.~Hayashi and T.~Shirafuji,
``Gravity from Poincar\'e gauge theory of the fundamental particles.
I. General formulation,''
Prog.\ Theor.\ Phys.\ \textbf{64}, 866 (1980).

\bibitem{Neville1980}
D.~E.~Neville,
``Gravity theories with propagating torsion,''
Phys.\ Rev.\ D~\textbf{21}, 867 (1980).

\bibitem{YoNester1999}
H.-J.~Yo and J.~M.~Nester,
``Hamiltonian analysis of Poincar\'e gauge theory,''
Int.\ J.\ Mod.\ Phys.\ D~\textbf{8}, 459 (1999).

\bibitem{MukhanovFeldmanBrandenberger1992}
V.~F.~Mukhanov, H.~A.~Feldman, and R.~H.~Brandenberger,
``Theory of cosmological perturbations,''
Phys.\ Rept.\ \textbf{215}, 203 (1992).

\bibitem{PeterPintoNeto2008}
P.~Peter and N.~Pinto-Neto,
``Primordial perturbations in a non-singular bouncing universe model,''
Phys.\ Rev.\ D~\textbf{78}, 063506 (2008).

\bibitem{Quintin2015}
J.~Quintin, Y.-F.~Cai, and R.~H.~Brandenberger,
``Matter creation in a nonsingular bouncing cosmology,''
Phys.\ Rev.\ D~\textbf{90}, 063507 (2015).

\bibitem{Cai2014}
Y.-F.~Cai, D.~A.~Easson, and R.~Brandenberger,
``Towards a nonsingular bouncing cosmology,''
JCAP \textbf{08}, 020 (2012).

\bibitem{Hehl1976}
F.~W.~Hehl, P.~von~der~Heyde, G.~D.~Kerlick, and J.~M.~Nester,
``General relativity with spin and torsion: Foundations and prospects,''
Rev.\ Mod.\ Phys.\ \textbf{48}, 393 (1976).

\bibitem{KaloPerPadilla2014}
N.~Kaloper and A.~Padilla,
``Sequestering the standard model vacuum energy,''
Phys.\ Rev.\ Lett.\ \textbf{112}, 091304 (2014).

\bibitem{Poplawski2016}
N.~J.~Pop\l{}awski,
``Universe in a black hole in Einstein-Cartan gravity,''
Astrophys.\ J.\ \textbf{832}, 96 (2016).

\bibitem{LISA2017}
LISA Collaboration,
``Laser Interferometer Space Antenna,''
arXiv:1702.00786 (2017).

\bibitem{ET2010}
M.~Punturo \etal,
``The Einstein Telescope: A third-generation gravitational wave
observatory,''
Class.\ Quant.\ Grav.\ \textbf{27}, 194002 (2010).

\bibitem{Saadeh2016}
D.~Saadeh \etal,
``How isotropic is the universe?''
Phys.\ Rev.\ Lett.\ \textbf{117}, 131302 (2016).

\bibitem{Legner2025}
C.~Legner, W.~J.~Handley, W.~E.~V.~Barker, and S.~S.~Ormondroyd,
``Torsion condensation in Poincar\'e gauge theory,''
arXiv:2507.09228 (2025).

\bibitem{Fabbri2025}
L.~Fabbri,
``Propagating torsion as an axial-vector field,''
arXiv:2502.17979 (2025).

\bibitem{Izaurieta2020}
F.~Izaurieta \etal,
``Torsion fermion interaction and the Hubble parameter,''
Phys.\ Rev.\ D~\textbf{101}, 064007 (2020).

\bibitem{Akhshabi2023}
S.~Akhshabi and E.~Zamani,
``Cosmological constraints on Einstein-Cartan torsion,''
Class.\ Quant.\ Grav.\ \textbf{40}, 065006 (2023).

\end{thebibliography}

\end{document}
